\documentclass[12pt,a4paper]{article}
\usepackage[utf8]{inputenc}
\usepackage[T1]{fontenc}
\usepackage[lithuanian]{babel}
\usepackage{lmodern}
\usepackage{amsmath, amssymb}
\usepackage{geometry}
\geometry{left=2cm, right=2cm, top=2.5cm, bottom=2.5cm}

\begin{document}

\begin{center}
    {\Large \textbf{Kontrolinis darbas: Daugianariai}}\\[0.3cm]
    {\large \textbf{Platono variantas -- Atsakymai ir sprendimai}}
\end{center}
\vspace{0.5cm}

\begin{enumerate}
    \item \textbf{Reiškinių tipai:}
    \begin{itemize}
        \item[a)] $\displaystyle \frac{x^4 - \sqrt{5}}{3}$ -- \textbf{R, S} (Racionalusis sveikas). 
        \item[b)] $\displaystyle \frac{4}{x^3 + 2}$ -- \textbf{R, T} (Racionalusis trupmeninis).
        \item[c)] $\displaystyle x^{-4} + \sqrt{2} x$ -- \textbf{R, T} (Racionalusis trupmeninis). $x^{-4} = \frac{1}{x^4}$.
        \item[d)] $\displaystyle \sqrt{x^2 - 9}$ -- \textbf{I} (Iracionalusis).
        \item[e)] $\displaystyle 2x + 3^{-2}$ -- \textbf{R, S} (Racionalusis sveikas). 
        \item[f)] $\displaystyle \sqrt{2}x - 7$ -- \textbf{R, S} (Racionalusis sveikas).
    \end{itemize}

    \vspace{0.3cm}
    \item \textbf{Vienanarių daugyba:} \\
    $$ \left( -2x^2y^{-3} \right)^5 \cdot \left( 2x^{-3}y^4 \right)^{-4} = (-2)^5 \cdot x^{10} \cdot y^{-15} \cdot 2^{-4} \cdot x^{12} \cdot y^{-16} $$
    $$ = -32 \cdot x^{10} \cdot y^{-15} \cdot \frac{1}{16} \cdot x^{12} \cdot y^{-16} = -2 \cdot x^{10+12} \cdot y^{-15-16} = \mathbf{-2x^{22}y^{-31}} \text{ arba } \mathbf{-\frac{2x^{22}}{y^{31}}} $$
    \textbf{Paaiškinimas:} Tai \textbf{nėra vienanaris}, nes $y$ laipsnio rodiklis yra neigiamas.

    \vspace{0.3cm}
    \item \textbf{Skaičiavimas be skaičiuotuvo:} \\
    Skaitiklis: $68^3 + 32^3 = (68 + 32)(68^2 - 68 \cdot 32 + 32^2) = 100 \cdot (68^2 - 68 \cdot 32 + 32^2)$. \\
    Vardiklis: $0,2^{-2} = \left(\frac{2}{10}\right)^{-2} = \left(\frac{1}{5}\right)^{-2} = 5^2 = 25$. \\
    Reiškinys skliaustuose: $68 \cdot 36 = 68 \cdot (68 - 32) = 68^2 - 68 \cdot 32$. \\
    Vardiklis tampa: $25 \cdot (68^2 - 68 \cdot 32 + 32^2)$. \\
    Padaliję gauname: $\frac{100}{25} = \mathbf{4}$.

    \vspace{0.3cm}
    \item \textbf{Formulių įrodymai:}
    \begin{itemize}
        \item[a)] $\frac{x^m}{x^n} = \frac{\overbrace{x \cdot x \cdot \dots \cdot x}^{m}}{\underbrace{x \cdot x \cdot \dots \cdot x}_{n}} = \underbrace{x \cdot x \cdot \dots \cdot x}_{m-n} = \mathbf{x^{m-n}}.$ (kur $m > n$).
        \item[b)] $(x-y)^3 = (x-y)(x-y)^2 = (x-y)(x^2 - 2xy + y^2) = x^3 - 2x^2y + xy^2 - x^2y + 2xy^2 - y^3 = \mathbf{x^3 - 3x^2y + 3xy^2 - y^3}.$
    \end{itemize}

    \vspace{0.3cm}
    \item \textbf{Didžiausia ir mažiausia reikšmės:} \\
    $-5x^2 + 2x + 1 = -5\left(x^2 - \frac{2}{5}x\right) + 1 = -5\left(x^2 - 2 \cdot x \cdot \frac{1}{5} + \frac{1}{25} - \frac{1}{25}\right) + 1$ \\
    $= -5\left(\left(x - \frac{1}{5}\right)^2 - \frac{1}{25}\right) + 1 = -5\left(x - \frac{1}{5}\right)^2 + \frac{1}{5} + 1 = \mathbf{-5\left(x - \frac{1}{5}\right)^2 + 1\frac{1}{5}}$ \\
    (Galima rašyti ir dešimtainėmis trupmenomis: $-5(x - 0,2)^2 + 1,2$). \\
    \textbf{Didžiausia reikšmė:} $\mathbf{1\frac{1}{5}}$ (arba $1,2$), kai $x = \frac{1}{5}$ (arba $0,2$). \\
    \textbf{Mažiausios nėra} (artėja į $-\infty$).

    \vspace{0.3cm}
    \item \textbf{Sudėtingesnis reiškinio pertvarkymas:} \\
    Iškeliame $(x+1)$: $(x+1)\left[ (x+1)^2 - (x+2)^2 \right]$ \\
    $(x+1) [ (x^2 + 2x + 1) - (x^2 + 4x + 4) ] = (x+1)[ -2x - 3 ] = -2x^2 - 3x - 2x - 3 = -2x^2 - 5x - 3$. \\
    Išskiriame pilnąjį kvadratą: $-2\left(x^2 + \frac{5}{2}x\right) - 3 = -2\left(x^2 + 2 \cdot \frac{5}{4}x + \frac{25}{16} - \frac{25}{16}\right) - 3 = -2\left(x + \frac{5}{4}\right)^2 + \frac{25}{8} - \frac{24}{8} = \mathbf{-2\left(x + \frac{5}{4}\right)^2 + \frac{1}{8}}$. \\
    Gavome koeficientus: $a = -2$, $b = \frac{5}{4}$, $c = \frac{1}{8}$. \\
    Skaičiuojame reiškinio po šaknimi reikšmę: \\
    $\sqrt{-2a - 4b + 8c + 7} = \sqrt{-2(-2) - 4\left(\frac{5}{4}\right) + 8\left(\frac{1}{8}\right) + 9} = \sqrt{4 - 5 + 1 + 9} = \mathbf{\sqrt{9}} = \mathbf{3}$.

    \vspace{0.3cm}
    \item \textbf{Trijų kintamųjų mažiausia reikšmė:} \\
    $(x^2 + 6x + 9) - 9 + (y^2 - 4y + 4) - 4 + (z^2 - 8z + 16) - 16$ \\
    $= (x+3)^2 + (y-2)^2 + (z-4)^2 - 29$ \\
    Kadangi pilnieji kvadratai visada neneigiami, \textbf{mažiausia reikšmė yra $\mathbf{-29}$} (kai $x=-3, y=2, z=4$).

\end{enumerate}

\end{document}
